\section{Maximum GCD on Whiteboard}
\label{sec:n-whiteboard}
\problemlink{https://codeforces.com/problemset/problem/2156/C}{Codeforces 2156C}
\source{Codeforces 2156C (Notion: Number Theory, 1400)}{Medium--Hard}

\problemstatement{$n\le2\cdot10^5$ numbers $1\le a_i\le n$. You may erase at
most $k$ numbers, and any number of times split $x\ge3$ into
$x_1\le x_2\le x_3$ with $x_1+x_2+x_3=x$, keeping only $x_1$ and $x_3$.
Maximise the gcd of the final board.}

\begin{patternbox}
Fix the answer $d$ and count what \emph{cannot} be fixed for it; values are
$\le n$, so every candidate $d$ can be checked with prefix counts in
$O(n/d)$.
\end{patternbox}

\subsection*{Approach and intuition}

Fix $d$. Multiples of $d$ are already fine. A non-multiple $x$ must be split
into pieces that are multiples of $d$ (or be erased). With $x_1=d$,
$x_3=td$ and $d\le x_2=x-d-td\le td$ we can reach every $x$ with
$d(t+2)\le x\le d(2t+1)$ for some $t\ge1$; these intervals cover $\{3d\}$ and
all of $[4d,\infty)$. Smaller $x$ (below $4d$, not a multiple) cannot work:
$x_1\le x/3<d$ and pieces never grow. (Verified by exhaustive search for
$d\le6$, $x<60$.)

\begin{ideabox}[Criterion]
$d$ is achievable iff $\#\{i : a_i<4d,\ d\nmid a_i\}\le k$. Take the largest
such $d$ ($d=1$ always works).
\end{ideabox}

With a prefix count array, the number of values $<4d$ is $pre[\min(4d-1,n)]$
and the good ones among them are the counts of $d,2d,3d$.

\subsection*{Algorithm}
\begin{algorithmic}[1]
\State $cnt[v]$, $pre[v]$ over $1..n$
\For{$d=n$ \textbf{down to} $1$}
  \State $bad\gets pre[\min(4d-1,n)]-cnt[d]-cnt[2d]-cnt[3d]$ (existing indices)
  \If{$bad\le k$} \Return $d$ \EndIf
\EndFor
\end{algorithmic}

\dryrun{$[4,9,6,8,2,6,7,8,2]$, $k=1$: $d=2$: values $<8$ that are odd: only
$7$, so $bad=1\le1$; larger $d$ fail. Answer $2$. With a second $7$ added,
$bad=2>1$ and the answer drops to $1$. \checkmark}

\subsection*{C++17 implementation}
\cppfile{number-theory/17_maximum_gcd_on_whiteboard.cpp}

\begin{complexitybox}
\textbf{Time:} $O(n)$ per test. \textbf{Space:} $O(n)$.
\end{complexitybox}

\begin{pitfallbox}
\begin{itemize}
  \item The threshold is $4d$, not $3d$: $x=7$, $d=2$ cannot be split
        ($2+2+3$ leaves $3$).
  \item Erase removes one copy only; count with multiplicity.
\end{itemize}
\end{pitfallbox}
